\SourceLeaf{153}keitsrechnung ganz unabhängige Antwort gegeben werden soll. Nur kann ich nicht wohl einräumen, das, was man auf diese Art erhält, den \textls{wahrscheinlichsten} Werth zu nennen. In der That ist die Aufgabe, den wahrscheinlichsten Werth zu finden, eine mathematisch ganz bestimmte, die aber ihrer Natur nach die Kenntniss des Fehlergesetzes voraussetzt und \textls{nur} in dem einzigen Falle, wo diess durch die Form $e^{-kkxx}$ ausgedrückt wird, auf die arithmetischen Mittel führt. Allgemein zu reden ist die Wahrscheinlichkeit des wahrscheinlichsten Werths auch nur unendlich klein, nemlich wenn $a$ der wahrscheinlichste Werth aus einer \textls{stetigen} Gesammtheit ist, so bedeutet diess im Grunde nur so viel, dass

\Needspace{7\baselineskip}
\begin{quote}
die Wahrscheinlichkeit, der wahre Werth liege zwischen $a-\omega$ und $a+\omega$, grösser ist als die Wahrscheinlichkeit, der wahre Werth liege zwischen irgend einem andern Paar eben so weiter Grenzen, in so fern $\omega$ unendlich klein ist.
\end{quote}

Genau besehen hat aber eben deshalb solcher wahrscheinlichster Werth nur wenig praktisches Interesse, viel weniger als derjenige Werth, wobei der zu befürchtende Irrthum im Durchschnitt am wenigsten schädlich ist, daher ich (ausser andern freilich eben so wichtigen oder noch viel wichtigern Gründen) dieses zweite mit dem ersten ja nicht zu verwechselnde Princip vorgezogen habe.

. . . . . . . . . . . .

\Needspace{6\baselineskip}
\begin{center}\textsc{Gauss} an \textsc{Bessel}. Göttingen, 28. Februar 1839.\end{center}

. . . . . . . . . . . . . . . .

Ihren Aufsatz in den Astronomischen Nachrichten über die Annäherung des Gesetzes für die Wahrscheinlichkeit aus zusammengesetzten Quellen entspringender Beobachtungsfehler an die Formel $e^{-\frac{xx}{hh}}$ habe ich mit grossem Interesse gelesen; doch bezog sich, wenn ich aufrichtig sprechen soll, dieses Interesse weniger auf die Sache selbst, als auf Ihre Darstellung. Denn jene ist mir seit vielen Jahren familiär, während ich selbst niemals dazu gekommen bin, die Entwickelung vollständig auszuführen.

